\chapter{Fractions and Rational Expressions}

Simplify completely; combine into a single fraction where applicable.

\section{Core Rules}

\begin{rulebox}{Factor first, then cancel --- and cancel only factors}
A fraction may be reduced only by a factor common to the \emph{whole}
numerator and the \emph{whole} denominator:
\[
  \frac{x^{3}+x^{2}}{x^{2}+x}=\frac{x^{2}(x+1)}{x(x+1)}=x .
\]
What is never allowed is striking out a term of a sum:
\[
  \frac{x+3}{x}\neq 3,\qquad
  \frac{a+b}{a+c}\neq\frac{b}{c} .
\]
If the numerator or denominator is a sum, factor it before cancelling
anything. If it will not factor, nothing cancels.
\end{rulebox}

\begin{rulebox}{Adding: one denominator, built from the factorisations}
The least common denominator is the product of all distinct factors, each
taken to its highest occurring power. Factor both denominators first:
\[
  \frac{4}{x^{2}-2x}+\frac{2}{x}
  =\frac{4}{x(x-2)}+\frac{2}{x}
  =\frac{4+2(x-2)}{x(x-2)}
  =\frac{2x}{x(x-2)}
  =\frac{2}{x-2}.
\]
Multiplying denominators blindly gives a common denominator too, just a
needlessly large one.
\end{rulebox}

\begin{rulebox}{Dividing and complex fractions}
Division is multiplication by the reciprocal,
$\dfrac{A}{B}\div\dfrac{C}{D}=\dfrac{A}{B}\cdot\dfrac{D}{C}$. For a fraction
whose parts are themselves fractions, either simplify the top and the bottom
to a single fraction each and then divide, or multiply top and bottom by the
overall common denominator. Both work; pick one and use it consistently.
\end{rulebox}

\begin{rulebox}{Restrictions travel with the expression}
$\dfrac{x^{2}-1}{x-1}=x+1$ holds \textbf{for $x\neq1$}. Cancelling a factor
silently deletes a point from the domain, and the deleted point is often
exactly the one the question is about --- a difference quotient at $h=0$, a
demand function at a choke price, a limit at a removable discontinuity.
State the excluded values.
\end{rulebox}

\section{Worked Examples}

\begin{worked}{(a) Simplify the complex fraction $\dfrac{\dfrac{1}{x-1}-\dfrac{1}{x+1}}{\dfrac{1}{x-1}+\dfrac{1}{x+1}}$}
\wstep{Step 1. Note the restrictions before simplifying.}
The inner fractions need $x\neq1$ and $x\neq-1$; the outer division needs a
non-zero denominator, which will turn out to exclude $x=0$ as well. Carry
these along.

\wstep{Step 2. Combine the numerator over the common denominator.}
\[
  \frac{1}{x-1}-\frac{1}{x+1}
  =\frac{(x+1)-(x-1)}{(x-1)(x+1)}
  =\frac{2}{x^{2}-1}.
\]
\wtrap{$(x+1)-(x-1)=2$, not $2x$. When subtracting, bracket the second
numerator before distributing the minus sign; a dropped bracket here is the
single most common error in the chapter.}

\wstep{Step 3. Combine the denominator the same way.}
\[
  \frac{1}{x-1}+\frac{1}{x+1}
  =\frac{(x+1)+(x-1)}{(x-1)(x+1)}
  =\frac{2x}{x^{2}-1}.
\]

\wstep{Step 4. Divide --- and watch the common denominator cancel.}
\[
  \frac{2/(x^{2}-1)}{2x/(x^{2}-1)}
  =\frac{2}{x^{2}-1}\cdot\frac{x^{2}-1}{2x}
  =\frac{2}{2x}
  =\boxed{\dfrac{1}{x}},\qquad x\neq0,\pm1 .
\]
The whole point of Steps~2 and~3 was to reach a single fraction over a
single fraction, where the shared denominator $x^{2}-1$ disappears in one
stroke.
\end{worked}

\begin{worked}{(b) The difference quotient of $f(x)=\dfrac1x$}
\wstep{Step 1. Write out $f(x+h)$ correctly.}
$f(x+h)=\dfrac{1}{x+h}$ --- substitute $x+h$ wherever $x$ appeared. It is
\emph{not} $\dfrac1x+h$.

\wstep{Step 2. Assemble the quotient.}
\[
  \frac{f(x+h)-f(x)}{h}=\frac{\dfrac{1}{x+h}-\dfrac1x}{h}.
\]
Setting $h=0$ now gives $\tfrac00$, which says nothing. The algebra exists
precisely to remove that.

\wstep{Step 3. Combine the top over the common denominator $x(x+h)$.}
\[
  \frac{1}{x+h}-\frac1x=\frac{x-(x+h)}{x(x+h)}=\frac{-h}{x(x+h)} .
\]
The numerator now carries an explicit factor $h$ --- which is the entire
purpose of the manipulation.

\wstep{Step 4. Divide by $h$, that is, multiply by $1/h$.}
\[
  \frac{-h}{x(x+h)}\cdot\frac1h=\frac{-1}{x(x+h)} .
\]
The cancellation is legal for every $h\neq0$, and $h\neq0$ is exactly the
case the difference quotient concerns.

\wstep{Step 5. Read off the limit.}
Letting $h\to0$ gives $-1/x^{2}$, the derivative of $1/x$. No limit rule was
needed; the algebra did all the work, and this is why the chapter comes
before the differentiation volume.
\wtrap{Do not cancel $h$ before the numerator has been written as a single
fraction. In $\bigl(\tfrac{1}{x+h}-\tfrac1x\bigr)/h$ there is no factor $h$
to cancel yet --- it only appears after Step~3.}
\end{worked}

\section{Practice Problems}

\begin{problems}
  \item $\displaystyle\frac{3a^{2}b}{c}\cdot\frac{c^{3}}{6a^{4}b^{2}}$
  \item $\displaystyle\frac{2}{x-3}-\frac{1}{x+3}$
  \item $\displaystyle\frac{x}{x+1}+\frac{1}{x+1}$
  \item $\displaystyle\frac{a/b+b/a}{a/b-b/a}$
  \item $\displaystyle\frac{x^{3}+x^{2}}{x^{2}+x}$
  \item Difference quotient: $\displaystyle\frac{f(x+h)-f(x)}{h}$
        for $f(x)=x^{3}$.
  \item Difference quotient: $\displaystyle\frac{f(x+h)-f(x)}{h}$
        for $f(x)=\dfrac{1}{x}$.
  \item $\displaystyle\frac{x^{4}-1}{x^{2}-1}$
  \item $\displaystyle\frac{4}{x^{2}-2x}+\frac{2}{x}$
  \item $\displaystyle\frac{1+x+x^{2}}{1-x^{3}}$
  \item $\displaystyle\frac{a^{3}+b^{3}}{a+b}$
  \item $\displaystyle\frac{x^{2}-2x+1}{x^{2}-1}$
  \item $\displaystyle\frac{\dfrac{1}{x-1}-\dfrac{1}{x+1}}
        {\dfrac{1}{x-1}+\dfrac{1}{x+1}}$
  \item MRS for Cobb--Douglas $U=x^{\alpha}y^{1-\alpha}$: simplify
        $\dfrac{\partial U/\partial x}{\partial U/\partial y}$.
  \item $\displaystyle\frac{2a^{2}-8}{a^{2}-4a+4}$
  \item Express $1+\dfrac{1}{x}$ as a single fraction.
  \item $\displaystyle\frac{x^{2}+6x+9}{(x+3)^{2}}$
  \item $\displaystyle\frac{x^{2}+4x+3}{x^{2}+5x+6}$
  \item $\displaystyle\frac{1}{2}\cdot\frac{x+1}{x+2}\div\frac{x+1}{2x+4}$
  \item $\displaystyle\frac{x^{3}-27}{x-3}$
  \item $\displaystyle\frac{ab+a+b+1}{a+1}$
  \item Solve for $x$: $\displaystyle\frac{1}{x}+\frac{1}{x+1}=\frac{5}{6}$
  \item $\displaystyle\frac{n!}{(n-2)!}$
  \item $\displaystyle\frac{a-b}{b-a}$
  \item $\displaystyle\frac{x^{2}+xy}{x^{2}-y^{2}}$
  \item $\displaystyle\frac{4x+8}{x^{2}+2x}$
  \item $\displaystyle\frac{1}{x(x-1)}-\frac{1}{x(x+1)}$
  \item $\displaystyle\frac{(x-2)^{3}}{(x-2)^{2}(x+1)}$
  \item Solve for $x$: $\displaystyle\frac{x+1}{x-1}=3$
  \item Solve for $x$: $\displaystyle\frac{2}{x}=\frac{3}{x+1}$
\end{problems}

\section{Solutions}

\begin{solutions}
  \item $\displaystyle\frac{c^{2}}{2a^{2}b}$
  \item $\displaystyle\frac{x+9}{x^{2}-9}$
  \item $1$
  \item $\displaystyle\frac{a^{2}+b^{2}}{a^{2}-b^{2}}$
  \item $x$
  \item $3x^{2}+3xh+h^{2}$
  \item $\displaystyle-\frac{1}{x(x+h)}$
  \item $x^{2}+1$
  \item $\displaystyle\frac{2}{x-2}$
  \item $\displaystyle\frac{1}{1-x}$
  \item $a^{2}-ab+b^{2}$
  \item $\displaystyle\frac{x-1}{x+1}$
  \item $\displaystyle\frac{1}{x}$
  \item $\displaystyle\frac{\alpha y}{(1-\alpha)x}$
  \item $\displaystyle\frac{2(a+2)}{a-2}$
  \item $\displaystyle\frac{x+1}{x}$
  \item $1$
  \item $\displaystyle\frac{x+1}{x+2}$
  \item $1$
  \item $x^{2}+3x+9$
  \item $b+1$
  \item $5x^{2}-7x-6=0\Rightarrow x=2$ (or $-\tfrac35$, rejected if
        positivity required)
  \item $n(n-1)$
  \item $-1$
  \item $\displaystyle\frac{x}{x-y}$
  \item $\displaystyle\frac{4}{x}$
  \item $\displaystyle\frac{2}{x(x^{2}-1)}$
  \item $\displaystyle\frac{x-2}{x+1}$
  \item $x=2$
  \item $x=2$
\end{solutions}
